\section{Objective}\label{sec:objective}
Candidates are ranked by the veto-free continuous objective \code{W\_ERGOLIGHT\_v1} (draft~d, requirement set
\code{RS-karambit-v3}; \src{crates/g3/src/ergolight.rs}{\code{ergolight.rs}}):
\begin{equation}
\mathrm{score}_{100} = 100\,E \sum_{i=1}^{21} w_i\,s_i,
\qquad
E = \prod_{k=1}^{8} e_k^{\beta_k},
\end{equation}
where $s_i \in [0,1]$ are normalized physical term scores with weights $\sum_{i=1}^{21} w_i = 1$, and
$E \in [0.3, 1.0]$ is a geometric-mean ergonomic multiplier over eight hand-zone compatibility factors
$e_k \in [0.3, 1.0]$ with exponents $\sum_{k=1}^{8} \beta_k = 1$. Because $E$ multiplies the entire structural
and mass sum rather than contributing an additive bonus, a candidate cannot trade hand comfort for grams: any
violation of the anatomical target bands scales down every point earned elsewhere. Conversely, once all eight hand-zone
measurements lie within their target plateau bands, $e_k = 1$ and $E = 1$, so ergonomics never distorts mass or
strength trade-offs among already comfortable knives.

\begin{table}[t]
\caption{Weights $w_i$ and ergonomic exponents $\beta_k$ of \code{W\_ERGOLIGHT\_v1} (draft~d)}\label{tab:weights}
\centering
\setlength{\tabcolsep}{3.2pt}
\begin{tabular}{@{}lrlr@{}}
\toprule
Additive term $s_i$ & $w_i$ & Ergonomic factor $e_k$ & $\beta_k$ \\
\midrule
Mass (linear to 25~g) & 0.3700 & Ring lip (\code{lip\_w}, bore) & 0.18 \\
Solid tip 3D (TET10) & 0.0700 & Grip section (perim., $H/T$) & 0.18 \\
Gate-S edge sharpness & 0.0500 & Hand-contact edge radii & 0.18 \\
Gate-S point profile & 0.0500 & Finger groove indexing & 0.15 \\
AM build \& drainage & 0.0500 & CoM \& spin inertia & 0.10 \\
CAD watertight census & 0.0400 & Grip chord length $L_g$ & 0.08 \\
Karambit ring \& claw & 0.0400 & Guard height $h_{\text{guard}}$ & 0.08 \\
Form fairness & 0.0400 & Grip texture amp./crest & 0.05 \\
Hiker lateral (180~N) & 0.0305 & & \\
Hunter lateral (260~N) & 0.0305 & & \\
Yield sequence ($\ge 1.2$) & 0.0305 & & \\
Ring back (30~N, SF~2) & 0.0261 & & \\
Ring bearing (300~N, SF~2) & 0.0218 & & \\
Blade body twist & 0.0218 & & \\
Blade tip twist & 0.0218 & & \\
3-level plate GCI & 0.0200 & & \\
Strict lateral (260~N) & 0.0174 & & \\
Blade buckling $\lambda$ & 0.0174 & & \\
Human-max twist/spine & 0.0174 & & \\
Goodman L7 fatigue & 0.0174 & & \\
Hollow grip $\mathrm{SF}_{\text{grip}}$ & 0.0174 & & \\
\bottomrule
\end{tabular}
\end{table}

Table~\ref{tab:weights} lists the 21 additive weights $w_i$ and 8 ergonomic exponents $\beta_k$. Mass carries the single
largest weight ($w_{\text{mass}} = 0.370$), mapped linearly from $0$ at $160$~g to $1$ at $25$~g so lighter designs are
always rewarded. In draft~d, the 3D solid-tip pry term carries $w_{\text{tip}} = 0.070$ to drive tip-root reinforcement
under the 1.82~N$\cdot$m tip pry, while Gate-S edge sharpness, Gate-S point profile, and AM buildability carry $0.050$
each. Every structural safety-factor ratio $v = \mathrm{SF}/\mathrm{SF}_{\text{req}}$ is mapped through a smooth
knee function (\code{sfk} in \src{crates/g3/src/ergolight.rs}{\code{ergolight.rs}}),
\begin{equation}
\operatorname{sfk}(v) = 0.94\,\min(\max(v, 0), 1) + 0.06\,\bigl(1 - e^{-\max(0, v - 1)/0.15}\bigr),
\end{equation}
which awards 94\,\% of the term credit continuously up to first-yield adequacy ($v = 1$) and saturates the remaining
6\,\% exponentially above $v = 1$ to reward margin against mesh and manufacturing variation without encouraging
over-built mass. No term acts as a hard reject gate during evolution, preserving smooth fitness gradients across the
archive.
